% New English translation from checked cote 25, site 64 (PDF page 65).
% AI-translated by GPT-6 Astra (OpenAI, ChatGPT Pro mode).
\scanpage{25/64}
\mspage{\unsure{57}}

point rational over k.

A point of $\underline{\mathrm{Grass}}_n(\underline{E})$ with values in S, i.e. a loc free quotient
Module of rank n \ $\underline{G}$ of $\underline{E}$, canonically defines a subprescheme
of $P(\underline{E})$, namely $P(\underline{G})$. Such a subprescheme\var{\inserted{(but without prescribing the rank of $\underline{G}$)}}{(but without prescribing the rank of $\underline{G}$)}{insertion}{ega25_25_064-v0001} is called a \uline{linear}
% NOTE: The last letter of the French sous-variété is cut by the right edge of the scan.
\uline{subvariety} of $P(\underline{E})$ (relative to S when confusion is possible).
It is therefore a projective bundle, of relative dimension n-1 if $n\geqslant 1$,
(and empty if n=0). We verify immediately that the section of $\underline{\mathrm{Grass}}_n$
$(\underline{E})$, i.e. G, \ is known when we know the corresponding linear subvariety
of $P(\underline{E})$. \unsure{In} this way, \ the Grassmannian can
be interpreted as representing the functor “linear subvarieties
of relative dimension n-1 \var{\xout{\unsure{(resp. empty)}}}{}{typewriter strike-out}{ega25_25_064-v0002} of $P_{S'}$” for variable S'
\var{\xout{x}}{}{typewriter strike-out}{ega25_25_064-v0003}\var{\inserted[typed]{$(n\geqslant 1)$}}{$(n\geqslant 1)$}{insertion}{ega25_25_064-v0004}. It is moreover possible to give an intrinsic characterization
of this last functor, i.e. of the notion of a linear subvariety
\ of relative dimension m: these are the closed subpreschemes
of P, smooth over S and everywhere of relative dimension m, and of
“projective degree 1” \ at every $s\in S$; \ this characterization will
be given in a later Chapter, and we have no need of it here.

Again assuming $\underline{E}$ \ loc free\var{\inserted[typed]{of rank r}}{ of rank r}{insertion}{ega25_25_064-v0005}, let $\underline{E}$ \ be its dual. Then
by polarity we find a canonical isomorphism
\[ \underline{\mathrm{Grass}}_n(E) \hand{\simeq} \underline{\mathrm{Grass}}_{r-n}(\overset{\hand{\vee}}{\underline{E}}) \]
associating \ to the quotient \ $\underline{G}$ of $\underline{E}$ \ the quotient \ $\overset{\hand{\vee}}{\underline{E}}/\overset{\hand{\vee}}{\underline{G}}$ \ of $\overset{\hand{\vee}}{\underline{E}}$.
From the viewpoint of linear varieties, the linear variety L of rel dim
m \ in P corresponds to the linear variety\var{\inserted[typed]{, its dual $L^{o}$,}}{, its dual $L^{o}$,}{insertion}{ega25_25_064-v0006} of relative dim
\var{\xout{\unsure{(r-1)-m}}}{}{typewriter strike-out}{ega25_25_064-v0007} \ (r-1)-1 -m \ \ of $\overset{\hand{\vee}}{P}$,\var{\inserted{i.e. of relative codim m+1 in $\check{P}$}}{i.e. of relative codim m+1 in $\check{P}$}{insertion}{ega25_25_064-v0008} (NB \ r-1 \ is here the common relative dim
of \ P and $\overset{\hand{\vee}}{P}$ over S), which can be visualized geometrically as follows.
First putting \ n=r-1, we find an isomorphism
\[ \underline{P}(\overset{\hand{\vee}}{\underline{E}}) \hand{\simeq} \underline{\mathrm{Grass}}_{r-1}(\underline{E}) \]
\tnote{25-064-1}{The source's e/g uncertainty in De cette façon is marked on In this way. The symbol introducing the dual in the lower paragraph is E without a dual accent in the checked transcription and is retained in that position. The inserted dual description is punctuated as an apposition in English, keeping its variant and formula together.}

\mspage{64}
